\chapter{Euler--Riemann Factorization of the Canonical Negative Inverse}
\label{chap:euler-riemann-negative-inversion}

\status{Exact operator factorization and fixed-point geometry proved. The chapter does not assert that non-trivial xi zeros are fixed points of the negative inversion and makes no RH claim.}

Chapter~\ref{chap:li-negative-inverse} established the canonical projective coordinate
\[
 \Omega(s)=\frac{s}{1-s}
\]
and Li's coordinate
\[
 z_L(s)=1-\frac1s,
\]
with the exact identity
\begin{equation}
 \boxed{z_L(s)=-\frac1{\Omega(s)}}.
 \label{eq:g012-li-negative-inverse}
\end{equation}
The purpose of this chapter is not to redefine that map.  It factorizes the same
negative inverse into the holomorphic Riemann functional reflection and the
Euler half-turn, then transports the resulting operator algebra through the
centered coordinates used in Chapters~32--36.

\section{Primitive involutions in the projective coordinate}

Set
\[
 u=\Omega(s)=\frac{s}{1-s}.
\]
The holomorphic functional reflection
\[
 \mathcal R_s(s)=1-s
\]
becomes
\begin{equation}
 \boxed{\mathcal R_u(u)=\frac1u}.
 \label{eq:g012-Ru}
\end{equation}
Indeed,
\[
 \Omega(1-s)=\frac{1-s}{s}=\frac1{\Omega(s)}.
\]

The Euler half-turn is multiplication by the phase
\[
 e^{i\pi}=-1,
\]
so define
\begin{equation}
 \boxed{\mathcal E_u(u)=-u=e^{i\pi}u}.
 \label{eq:g012-Eu}
\end{equation}
Both maps are involutions on their natural projective domain.

\section{Factorization of the canonical negative inverse}

The two primitive involutions commute:
\[
 \mathcal R_u\mathcal E_u(u)
 =\frac1{-u}
 =-\frac1u
 =\mathcal E_u\mathcal R_u(u).
\]
Define
\[
 \mathcal N:=\mathcal R\circ\mathcal E
 =\mathcal E\circ\mathcal R.
\]
Then
\begin{equation}
 \boxed{\mathcal N_u(u)=-\frac1u=\frac{e^{i\pi}}u}.
 \label{eq:g012-Nu}
\end{equation}
Combining \cref{eq:g012-li-negative-inverse,eq:g012-Nu} gives the exact
crosswalk to the Chapter~25 canon:
\begin{equation}
 \boxed{
 z_L(s)
 =\mathcal N_u(\Omega(s))
 =\mathcal E_u\!\left(\mathcal R_u(\Omega(s))\right)
 =\mathcal R_u\!\left(\mathcal E_u(\Omega(s))\right).
 }
 \label{eq:g012-li-factorization}
\end{equation}
Thus SOH-L017 is retained unchanged, while SOH-G012 supplies an operator
factorization of the same rational map.

\begin{theorem}[Euler--Riemann negative-inversion algebra]
On the common projective domain,
\[
 \mathcal R^2=\mathcal E^2=\mathcal N^2=\mathrm{id},
 \qquad
 \mathcal R\mathcal E=\mathcal E\mathcal R=\mathcal N.
\]
Consequently
\[
 \boxed{\{\mathrm{id},\mathcal R,\mathcal E,\mathcal N\}\cong V_4},
\]
the Klein four group.
\end{theorem}

\begin{proof}
The involutivity of \(\mathcal R\) and \(\mathcal E\) follows immediately from
\(1/(1/u)=u\) and \(-(-u)=u\).  Equation~\eqref{eq:g012-Nu} gives
\(\mathcal N^2(u)=-1/(-1/u)=u\), while the displayed composition identity gives
commutativity and closure.  The four maps are distinct on a generic projective
point, so the resulting group is the Klein four group.
\end{proof}

\section{Centered conjugacies}

Use the centered coordinate
\[
 t=2s-1=\frac{u-1}{u+1}.
\]
Direct substitution gives
\begin{equation}
 \boxed{\mathcal R_t(t)=-t},
 \qquad
 \boxed{\mathcal E_t(t)=\frac1t},
 \qquad
 \boxed{\mathcal N_t(t)=-\frac1t}.
 \label{eq:g012-t-actions}
\end{equation}
The negative inversion in the centered chart is therefore exactly the product
of the Riemann sign reflection and the Euler reciprocal action.

With
\[
 z=s-\frac12=\frac t2,
\]
one obtains
\begin{equation}
 \boxed{\mathcal R_z(z)=-z},
 \qquad
 \boxed{\mathcal E_z(z)=\frac1{4z}},
 \qquad
 \boxed{\mathcal N_z(z)=-\frac1{4z}}.
 \label{eq:g012-z-actions}
\end{equation}
Finally, for the even quotient
\[
 w=z^2,
\]
we have
\begin{equation}
 \boxed{\mathcal R_w(w)=w},
 \qquad
 \boxed{\mathcal E_w(w)=\mathcal N_w(w)=\frac1{16w}}.
 \label{eq:g012-w-actions}
\end{equation}
Thus the square quotient removes the distinction between \(\mathcal E\) and
\(\mathcal N\) because it has already quotiented out the sign reflection
\(\mathcal R\).

\section{Fixed points and the critical line}

The affine fixed-point equation for the canonical negative inversion is
\[
 u=-\frac1u.
\]
Equivalently,
\begin{equation}
 \boxed{u^2=e^{i\pi}=-1}.
 \label{eq:g012-fixed-euler}
\end{equation}
Hence
\[
 \boxed{u_\pm=\pm i}.
\]
Mapping back with \(s=u/(1+u)\) gives
\begin{equation}
 \boxed{s_\pm=\frac12\pm\frac{i}{2}}.
 \label{eq:g012-fixed-s}
\end{equation}
Therefore
\[
 \boxed{\Re(s_\pm)=\frac12}.
\]

\begin{theorem}[Fixed pair of the canonical negative inversion]
The two affine fixed points of \(u\mapsto-1/u\) are \(u=\pm i\).  Their
projective preimages under \(u=\Omega(s)\) are
\[
 \boxed{s=\frac12\pm\frac{i}{2}},
\]
so both lie on the Riemann critical line.
\end{theorem}

\begin{warning}
This fixed-point statement is an operator-geometric theorem.  It does not state
that a non-trivial zero of \(\xi\) must be fixed by \(\mathcal N\), and it does
not imply the Riemann Hypothesis.
\end{warning}

In the centered coordinates the same fixed pair is
\[
 \boxed{t_\pm=\pm i},
 \qquad
 \boxed{z_\pm=\pm\frac{i}{2}},
 \qquad
 \boxed{w=-\frac14}.
\]

\section{Fixed-point stratification after the square quotient}

The quotient action
\[
 \mathcal N_w(w)=\frac1{16w}
\]
has the two fixed values
\begin{equation}
 \boxed{w=\pm\frac14}.
 \label{eq:g012-w-fixed}
\end{equation}
They have different origins.

For \(w=-1/4\), the preimages \(z=\pm i/2\) are genuine fixed points of
\(\mathcal N_z\):
\[
 \mathcal N_z\!\left(\pm\frac{i}{2}\right)
 =\pm\frac{i}{2}.
\]
For \(w=+1/4\), the preimages \(z=\pm1/2\) form a two-cycle:
\[
 \mathcal N_z\!\left(\frac12\right)=-\frac12,
 \qquad
 \mathcal N_z\!\left(-\frac12\right)=\frac12.
\]
The map \(w=z^2\) identifies that two-cycle, producing a second fixed value only
after quotienting.  Hence \(w=-1/4\) is the image of genuine negative-inversion
fixed points, whereas \(w=+1/4\) is a quotient-created fixed value.

\section{Logarithmic half-period form}

Let
\[
 u=e^\lambda.
\]
Modulo the full logarithmic period \(2\pi i\), the primitive maps lift to
\[
 \mathcal R_\lambda(\lambda)=-\lambda,
 \qquad
 \mathcal E_\lambda(\lambda)=\lambda+i\pi.
\]
Their composition is
\begin{equation}
 \boxed{\mathcal N_\lambda(\lambda)=i\pi-\lambda}.
 \label{eq:g012-log-N}
\end{equation}
The fixed lifts satisfy
\[
 \lambda=i\pi-\lambda+2\pi i k,
\]
so
\begin{equation}
 \boxed{\lambda_k=\frac{i\pi}{2}+\pi i k},
 \qquad k\in\mathbb Z.
 \label{eq:g012-log-fixed}
\end{equation}
Exponentiation collapses this family to the two classes \(u=+i\) and \(u=-i\).

\section{Proof firewall}

Proved in SOH-G012:
\begin{itemize}
 \item the Chapter~25 identity \(z_L(s)=-1/\Omega(s)\) is exactly
       \(\mathcal N_u(\Omega(s))\);
 \item \(\mathcal N=\mathcal R\mathcal E=\mathcal E\mathcal R\);
 \item \(\{\mathrm{id},\mathcal R,\mathcal E,\mathcal N\}\cong V_4\);
 \item the exact conjugacies in the \(t\), \(z\), and \(w\) coordinates;
 \item the affine fixed pair \(u=\pm i\) and its critical-line images
       \(s=1/2\pm i/2\);
 \item the distinct origins of the quotient fixed values \(w=\pm1/4\);
 \item the logarithmic lift \(\lambda\mapsto i\pi-\lambda\) and its fixed
       half-period family.
\end{itemize}

Not proved in SOH-G012:
\begin{itemize}
 \item that any non-trivial xi zero is fixed by \(\mathcal N\);
 \item that all xi zeros lie on the critical line;
 \item PF\(_\infty\);
 \item SOH-G003 real-rootedness;
 \item the Riemann Hypothesis.
\end{itemize}
